\chapter{Odwzorowania wieloliniowe}
\label{ch:lin8-wieloliniowe}
\index{odwzorowanie wieloliniowe@Odwzorowanie wieloliniowe}
\index{forma wieloliniowa@Forma wieloliniowa}
\index{forma symetryczna@Forma symetryczna}
\index{forma naprzemienna@Forma naprzemienna}

\noindent\textbf{Cel rozdziału.}
Odwzorowania liniowe zachowują kombinacje liniowe
jednego argumentu. W wielu konstrukcjach potrzebujemy
funkcji zależnych od kilku wektorów, liniowych
w każdym z nich oddzielnie. Takie odwzorowania
opisują formy dwuliniowe, wyznaczniki i formy
o wyższej liczbie argumentów.
Wyprowadzimy ich wzory współrzędnych oraz zbadamy
symetrię, naprzemienność i związane z nimi
przestrzenie liniowe. Nie będziemy jeszcze
zakładać istnienia iloczynów tensorowych;
ich konstrukcja jest tematem Rozdziału~\ref{ch:lin10-iloczyny-tensorowe}.

W całym rozdziale \(K\) jest ciałem,
\(0_K\ne1_K\). Wszystkie przestrzenie
wektorowe występujące w twierdzeniach
o wymiarze są skończenie wymiarowe,
chyba że wyraźnie zaznaczono inaczej.
Dopuszczamy wymiar \(0\) i dowolną
charakterystykę ciała.
Przyjmujemy \(\mathbb N_0=\{0,1,2,\ldots\}\);
pusta suma wynosi \(0_K\), a pusty iloczyn
skalarny wynosi \(1_K\).

\section{Definicja i podstawowe przykłady}
\label{sec:lin8-definicja}

\begin{definicja}[Odwzorowanie wieloliniowe]
\label{def:lin8-wieloliniowe}
Niech \(k\geq1\), a \(V_1,\ldots,V_k,W\)
będą przestrzeniami wektorowymi nad \(K\).
Odwzorowanie
\[
 F\colon V_1\times\cdots\times V_k\longrightarrow W
\]
nazywamy \emph{\(k\)-liniowym} lub
\emph{wieloliniowym}, jeżeli dla każdego
\(r\in\{1,\ldots,k\}\), dowolnych
\(v_i\in V_i\) przy \(i\ne r\),
\(u,v\in V_r\) oraz \(a,b\in K\) mamy
\[
 \begin{aligned}
 &F(v_1,\ldots,v_{r-1},au+bv,
                  v_{r+1},\ldots,v_k)\\
 &\qquad =
 aF(v_1,\ldots,v_{r-1},u,v_{r+1},\ldots,v_k)
 +bF(v_1,\ldots,v_{r-1},v,v_{r+1},\ldots,v_k).
 \end{aligned}
\]
Odwzorowanie o wartościach w \(K\)
nazywamy \emph{formą wieloliniową}.
Dla \(k=2\) używamy określenia
\emph{odwzorowanie dwuliniowe}.
Zbiór wszystkich \(k\)-liniowych
odwzorowań do \(W\) oznaczamy
przez \(\operatorname{Mult}_K(V_1,\ldots,V_k;W)\).
\end{definicja}

\begin{stwierdzenie}[Liniowość punktowa i argument zerowy]
\label{prop:lin8-punktowe}
Niech \(k\geq1\) oraz
\(F,G\in\operatorname{Mult}_K(V_1,\ldots,V_k;W)\).
Dla dowolnych \(a,b\in K\)
odwzorowanie \(aF+bG\), określone punktowo,
jest wieloliniowe. W szczególności
\(\operatorname{Mult}_K(V_1,\ldots,V_k;W)\)
jest przestrzenią wektorową nad \(K\).
Jeśli przynajmniej jeden argument
odwzorowania \(F\) jest wektorem zerowym,
to wartość \(F\) wynosi \(0_W\).
\end{stwierdzenie}
\begin{proof}
Dla ustalonych pozostałych argumentów
funkcje zależne od argumentu \(r\)
są liniowe. Ich kombinacja
\(aF+bG\) jest również liniowa
w tym argumencie.
Dotyczy to każdego \(r\), zatem
\(aF+bG\) jest wieloliniowe.
Pozostałe aksjomaty przestrzeni
wektorowej wynikają z punktowych
działań w \(W\).

Dla argumentu zerowego wykorzystujemy
liniowość w odpowiednim miejscu:
\(F(\ldots,0_{V_r},\ldots)
=F(\ldots,0_K\,v,\ldots)
=0_KF(\ldots,v,\ldots)=0_W\).
\end{proof}

\begin{przyklad}[Funkcjonały i iloczyny funkcjonałów]
\label{prz:lin8-iloczyn-funkcjonalow}
Dla \(k=1\) otrzymujemy zwykłe
odwzorowania liniowe; w szczególności
\(\operatorname{Mult}_K(V;K)=V^*\).
Jeżeli \(k\geq2\) i
\(\varphi_r\in V_r^*\)
dla \(r=1,\ldots,k\), to wzór
\[
 F(v_1,\ldots,v_k)
 :=\prod_{r=1}^k\varphi_r(v_r)
\]
określa formę wieloliniową.
Liniowość w argumencie \(r\)
wynika z liniowości \(\varphi_r\)
oraz rozdzielności mnożenia przez
pozostałe, ustalone czynniki.
Nie każda forma wieloliniowa musi być
jednym takim iloczynem.
\end{przyklad}

\begin{przyklad}[Mnożenie wielomianów]
\label{prz:lin8-wielomiany}
Niech \(V=K[t]\), przestrzeń
wielomianów jednej zmiennej
nad \(K\).
Odwzorowanie
\[
 M\colon V\times V\to V,\qquad
 M(p,q)=pq
\]
jest dwuliniowe dzięki rozdzielności
mnożenia wielomianów.
Nie jest natomiast liniowe jako
odwzorowanie z przestrzeni
\(V\times V\) wyposażonej w
działania punktowe: dla
\(p=q=1\) otrzymujemy
\(M((p,q)+(p,q))=4\cdot1\),
podczas gdy
\(M(p,q)+M(p,q)=2\cdot1\)
nad \(\mathbb R\).
Przykład odróżnia wieloliniowość
od zwykłej liniowości na produkcie.
\end{przyklad}

\begin{przyklad}[Wyznacznik jako forma wieloliniowa]
\label{prz:lin8-determinant}
Dla każdego \(n\geq1\)
odwzorowanie
\[
 D\colon (K^n)^n\to K,\qquad
 D(v_1,\ldots,v_n)
 =\det\begin{pmatrix}v_1&\cdots&v_n\end{pmatrix}
\]
jest \(n\)-liniowe.
Wynika to z liniowości wyznacznika
w każdej kolumnie, udowodnionej w
Stwierdzeniu~\ref{prop:lin6-transpozycja}.
Ponadto \(D\) znika, gdy dwa
argumenty są równe. Ta dodatkowa
własność będzie dla nas kluczowa.
\end{przyklad}

\begin{uwaga}[Umowa dla liczby argumentów zero]
\label{rem:lin8-zeroliniowe}
Dla \(k=0\) produkt pustej rodziny
przestrzeni jest zbiorem jednoelementowym
\(\{()\}\). Każde odwzorowanie
\(\{()\}\to W\) utożsamiamy z jego
jedyną wartością w \(W\).
Otrzymujemy naturalną identyfikację
\(\operatorname{Mult}_K(;W)\cong W\).
W szczególności przestrzeń
\(0\)-liniowych form skalarnych
jest równa \(K\).
Wszystkie późniejsze stwierdzenia
dotyczące sumowania po permutacjach
dla \(k=0\) interpretujemy z użyciem
jedynej permutacji pustego zbioru.
\end{uwaga}

\section{Współrzędne i wymiar przestrzeni odwzorowań}
\label{sec:lin8-wspolrzedne}

\begin{twierdzenie}[Rozwinięcie względem baz]
\label{thm:lin8-wspolrzedne}
Niech \(k\geq1\) i
\(V_1,\ldots,V_k,W\) będą
przestrzeniami skończenie wymiarowymi
o wymiarach odpowiednio
\(n_1,\ldots,n_k,m\).
Niech
\(B_r=(e^{(r)}_1,\ldots,e^{(r)}_{n_r})\)
będzie bazą \(V_r\).
Dla każdego układu wektorów
\[
 c_{i_1\ldots i_k}\in W,\qquad
 1\leq i_r\leq n_r,
\]
istnieje dokładnie jedno odwzorowanie
\(k\)-liniowe \(F\) spełniające
\[
 F(e^{(1)}_{i_1},\ldots,e^{(k)}_{i_k})
 =c_{i_1\ldots i_k}.
\]
Jeżeli
\(v_r=\sum_{i_r=1}^{n_r}
x^{(r)}_{i_r}e^{(r)}_{i_r}\),
to odwzorowanie to jest dane wzorem
\[
 F(v_1,\ldots,v_k)
 =\sum_{i_1=1}^{n_1}\cdots
  \sum_{i_k=1}^{n_k}
  \left(\prod_{r=1}^{k}x^{(r)}_{i_r}\right)
  c_{i_1\ldots i_k}.
\]
W szczególności
\[
 \dim\operatorname{Mult}_K(V_1,\ldots,V_k;W)
 =m\prod_{r=1}^{k}n_r.
\]
\end{twierdzenie}
\begin{proof}
Rozwińmy kolejno każdy argument
\(v_r\) względem bazy \(B_r\).
Wieloliniowość wymusza wskazaną
sumę, a więc odwzorowanie
o zadanych wartościach na
wszystkich krotkach bazowych
może być najwyżej jedno.

Odwrotnie, zadajmy dowolnie
\(c_{i_1\ldots i_k}\)
i \emph{zdefiniujmy} \(F\)
powyższym wzorem. Współrzędne
wektorów w ustalonych bazach
istnieją i są jednoznaczne.
Współrzędne kombinacji
\(av_r+bw_r\) są odpowiednimi
kombinacjami współrzędnych
\(v_r,w_r\); po podstawieniu
do wzoru rozdzielność mnożenia
dowodzi liniowości w argumencie
\(r\). Ponieważ \(r\) był dowolny,
\(F\) jest wieloliniowe.
Podstawienie wektorów bazowych
daje \(F(e^{(1)}_{i_1},\ldots,
e^{(k)}_{i_k})=c_{i_1\ldots i_k}\).

Odwzorowanie przypisujące \(F\)
rodzinę wszystkich wartości
\(c_{i_1\ldots i_k}\)
jest więc liniowym izomorfizmem
z przestrzenią
\(W^{\,n_1\cdots n_k}\).
Jej wymiar to
\(m\prod_r n_r\), co wynika
z wymiaru skończonej sumy prostej
(Twierdzenie~\ref{thm:lin2-wymiar-sumy-prostej}).
Jeśli któreś \(n_r=0\), dziedzina
produktu zawiera w każdym elemencie
argument zerowy, więc jedynym
odwzorowaniem wieloliniowym
jest zerowe. Zbiór rodzin
współczynników jest wtedy pusty,
a podany wymiar wynosi \(0\).
Jeśli \(m=0\), teza jest równie
bezpośrednia.
\end{proof}

\begin{wniosek}[Baza przestrzeni form wieloliniowych]
\label{cor:lin8-baza-wieloliniowych}
Jeżeli \(k\geq1\) oraz
\(V_r\) i \(B_r\) są jak w
Twierdzeniu~\ref{thm:lin8-wspolrzedne},
to formy
\[
 E_{i_1\ldots i_k}(v_1,\ldots,v_k)
 :=e^{(1),i_1}(v_1)\cdots
   e^{(k),i_k}(v_k),
\]
gdzie \(e^{(r),i}\) jest elementem
bazy dualnej do \(B_r\),
tworzą bazę przestrzeni
\(\operatorname{Mult}_K(V_1,\ldots,V_k;K)\).
W szczególności jej wymiar wynosi
\(\prod_{r=1}^{k}\dim V_r\).
\end{wniosek}
\begin{proof}
Na krotce wektorów bazowych
\(E_{i_1\ldots i_k}\) ma wartość
\(\prod_r\delta_{i_rj_r}\),
czyli \(1_K\) dokładnie wtedy,
gdy wszystkie indeksy się zgadzają,
a w pozostałych przypadkach \(0_K\).
Współrzędne z poprzedniego
twierdzenia pokazują, że każda
forma jest jednoznaczną kombinacją
tych funkcji. Ich liniowa
niezależność wynika także
bezpośrednio przez podstawianie
krotek bazowych.
Jeśli występuje przestrzeń zerowa,
rodzina form bazowych jest pusta,
zgodnie z wymiarem \(0\).
\end{proof}

\begin{przyklad}[Forma trójliniowa w \(\mathbb R^2\)]
\label{prz:lin8-trojliniowa}
Dla \(V=\mathbb R^2\) i
\(u=(u_1,u_2)\),
\(v=(v_1,v_2)\),
\(w=(w_1,w_2)\) niech
\[
 F(u,v,w)=u_1v_2w_1
          +2u_2v_1w_2.
\]
Jest to forma trójliniowa.
Jej jedyne niezerowe
współczynniki w bazie standardowej
to \(c_{121}=1\) i \(c_{212}=2\).
Wymiar przestrzeni wszystkich
form trójliniowych na
\((\mathbb R^2)^3\) wynosi
\(2^3=8\), mimo że tu
wykorzystujemy tylko dwa z ośmiu
współczynników.
Przykładowo
\(F((1,2),(3,4),(5,6))
=1\cdot4\cdot5+
 2\cdot2\cdot3\cdot6=92\).
\end{przyklad}

\begin{stwierdzenie}[Przekształcanie współczynników]
\label{prop:lin8-zmiana-baz}
Niech \(k\geq1\), a
\(F\colon V_1\times\cdots\times V_k\to W\)
będzie odwzorowaniem wieloliniowym
pomiędzy przestrzeniami
skończenie wymiarowymi.
Niech \(B_r=(e^{(r)}_i)\)
i \(C_r=(f^{(r)}_j)\)
będą uporządkowanymi bazami \(V_r\),
zaś \(D=(w_\alpha)\) i
\(E=(z_\beta)\) bazami \(W\).
Niech \(P_r\) spełnia
\([v]_{B_r}=P_r[v]_{C_r}\),
a \(Q\) spełnia
\([w]_D=Q[w]_E\).
Oznaczmy przez
\(a^\alpha_{i_1\ldots i_k}\)
współczynniki \(F\)
w bazach \(B_1,\ldots,B_k,D\),
a przez \(b^\beta_{j_1\ldots j_k}\)
współczynniki w bazach
\(C_1,\ldots,C_k,E\).
Wówczas
\[
 \begin{aligned}
 b^\beta_{j_1\ldots j_k}
 =\sum_{\alpha=1}^{\dim W}
  \ \sum_{i_1=1}^{\dim V_1}\cdots
  \sum_{i_k=1}^{\dim V_k}
 (Q^{-1})_{\beta\alpha}\,
 a^\alpha_{i_1\ldots i_k}
 \prod_{r=1}^{k}(P_r)_{i_rj_r}.
 \end{aligned}
\]
\end{stwierdzenie}
\begin{proof}
Dla każdego \(r\) z definicji
macierzy przejścia mamy
\(f^{(r)}_{j_r}
=\sum_{i_r}(P_r)_{i_rj_r}e^{(r)}_{i_r}\).
Podstawmy te rozwinięcia
do \(F(f^{(1)}_{j_1},\ldots,
f^{(k)}_{j_k})\) i kolejno
zastosujmy wieloliniowość.
Otrzymujemy sumę wartości
\(F(e^{(1)}_{i_1},\ldots,
e^{(k)}_{i_k})\)
z iloczynami odpowiednich
współczynników macierzy \(P_r\).

Wartość każdej z tych krotek
ma współrzędne
\(a^\alpha_{i_1\ldots i_k}\)
w bazie \(D\).
Zmiana współrzędnych wyjściowych
z \(D\) na \(E\) jest dana
macierzą \(Q^{-1}\).
Po podstawieniu i zebraniu
współczynnika przy \(z_\beta\)
dostajemy dokładnie wzór z tezy.
Jeżeli wymiar którejś przestrzeni
wynosi \(0\), obie strony są
pustymi rodzinami współczynników.
\end{proof}

\section{Formy dwuliniowe i ich macierze}
\label{sec:lin8-dwuliniowe}

\begin{definicja}[Forma dwuliniowa i macierz formy]
\label{def:lin8-forma-dwuliniowa}
Niech \(V\) będzie skończenie wymiarową
przestrzenią nad \(K\), z uporządkowaną
bazą \(B=(e_1,\ldots,e_n)\).
Dla formy dwuliniowej
\(b\colon V\times V\to K\)
jej \emph{macierzą w bazie \(B\)}
nazywamy
\[
 A_B(b):=(b(e_i,e_j))_{1\leq i,j\leq n}.
\]
Forma jest \emph{symetryczna},
jeżeli \(b(v,w)=b(w,v)\)
dla wszystkich \(v,w\in V\).
Jest \emph{naprzemienna},
jeżeli \(b(v,v)=0_K\)
dla każdego \(v\in V\).
Jest \emph{antysymetryczna},
jeżeli \(b(v,w)=-b(w,v)\)
dla wszystkich \(v,w\in V\).
\end{definicja}

\begin{twierdzenie}[Macierz i zmiana bazy formy dwuliniowej]
\label{thm:lin8-macierz-formy}
Niech \(b\) i \(B\) będą
jak w Definicji~\ref{def:lin8-forma-dwuliniowa},
a \(A=A_B(b)\).
Wtedy
\[
 b(v,w)=[v]_B^{\mathsf T}A[w]_B
 \qquad(v,w\in V).
\]
Każda macierz \(A\in
\operatorname{Mat}_{n\times n}(K)\)
określa w ten sposób dokładnie
jedną formę dwuliniową na \(V\).
Jeśli \(C\) jest inną bazą
i \([v]_B=P[v]_C\),
to
\[
 A_C(b)=P^{\mathsf T}A_B(b)P.
\]
Ponadto \(b\) jest symetryczna
dokładnie wtedy, gdy
\(A_B(b)^{\mathsf T}=A_B(b)\).
\end{twierdzenie}
\begin{proof}
Zapiszmy \(v=\sum_i x_ie_i\)
oraz \(w=\sum_j y_je_j\).
Z dwuliniowości wynika
\[
 b(v,w)=\sum_{i,j}x_iy_jb(e_i,e_j)
       =x^{\mathsf T}Ay.
\]
Zadając dowolną macierz \(A\)
i prawą stronę jako definicję \(b\),
otrzymujemy formę dwuliniową;
odwzorowanie jest liniowe oddzielnie
w \(x\) i \(y\), więc oddzielnie
w \(v\) i \(w\).
Jednoznaczność wynika
z wartości na wektorach bazy.

Przy zmianie bazy
\(x=P[v]_C\), \(y=P[w]_C\).
Stąd
\[
 b(v,w)=[v]_C^{\mathsf T}
          (P^{\mathsf T}AP)[w]_C.
\]
Jednoznaczność macierzy formy
daje żądaną równość.
Jeżeli \(b\) jest symetryczna,
to \(a_{ij}=b(e_i,e_j)
=b(e_j,e_i)=a_{ji}\).
Odwrotnie, gdy \(A=A^{\mathsf T}\),
otrzymujemy
\(x^{\mathsf T}Ay
=y^{\mathsf T}A^{\mathsf T}x
=y^{\mathsf T}Ax\),
więc \(b(v,w)=b(w,v)\).
W przypadku \(n=0\)
wszystkie tożsamości dotyczą
jedynych pustych macierzy.
\end{proof}

\begin{uwaga}[Kongruencja a podobieństwo]
\label{rem:lin8-kongruencja}
Prawo \(A_C=P^{\mathsf T}A_BP\)
dla form dwuliniowych jest inne
niż prawo \(P^{-1}A_BP\) dla
macierzy endomorfizmu; porównaj
Wniosek~\ref{cor:lin4-podobienstwo}.
W szczególności dwie macierze tej
samej formy nie muszą być podobne.
Na przykład nad \(\mathbb R\)
forma o macierzy \(I_2\)
w bazie standardowej ma
w bazie \(((2,0),(0,1))\)
macierz \(\operatorname{diag}(4,1)\);
obie macierze mają różne ślady,
więc nie są podobne.
\end{uwaga}

\begin{definicja}[Lewa i prawa podprzestrzeń radykalna]
\label{def:lin8-radykal}
Dla formy dwuliniowej
\(b\colon V\times V\to K\)
na przestrzeni skończenie wymiarowej
definiujemy jej \emph{lewą} i
\emph{prawą podprzestrzeń radykalną}:
\[
 \begin{aligned}
 \operatorname{Rad}_L(b)
 &:=\{v\in V:b(v,w)=0_K
                  \text{ dla każdego }w\in V\},\\
 \operatorname{Rad}_R(b)
 &:=\{w\in V:b(v,w)=0_K
                  \text{ dla każdego }v\in V\}.
 \end{aligned}
\]
Formę nazywamy \emph{niezdegenerowaną},
jeśli obie te podprzestrzenie
są zerowe.
\end{definicja}

\begin{twierdzenie}[Niezdegenerowanie i macierz]
\label{thm:lin8-niezdegenerowanie}
Niech \(\dim V=n<\infty\) i
\(b\colon V\times V\to K\)
będzie formą dwuliniową
o macierzy \(A\) w ustalonej bazie.
Wtedy współrzędne wektorów
w \(\operatorname{Rad}_L(b)\)
tworzą \(\ker A^{\mathsf T}\),
a współrzędne wektorów
w \(\operatorname{Rad}_R(b)\)
tworzą \(\ker A\).
W szczególności następujące
warunki są równoważne:
\[
 \begin{gathered}
 b\text{ jest niezdegenerowana},\\
 \operatorname{Rad}_L(b)=\{0\},\qquad
 \operatorname{Rad}_R(b)=\{0\},\\
 \det A\ne0_K.
 \end{gathered}
\]
\end{twierdzenie}
\begin{proof}
Dla \(x=[v]_B\), \(y=[w]_B\)
mamy \(b(v,w)=x^{\mathsf T}Ay\).
Wektor \(v\) należy do lewego
radykału wtedy i tylko wtedy, gdy
\(x^{\mathsf T}Ay=0_K\)
dla każdego \(y\in K^n\),
co jest równoważne
\(A^{\mathsf T}x=0\).
Analogicznie \(w\) należy
do prawego radykału dokładnie
wtedy, gdy \(Ay=0\).
Zatem wymiary obu radykałów
są równe \(n-\operatorname{rk}A\),
gdyż
\(\operatorname{rk}A^{\mathsf T}
=\operatorname{rk}A\)
na mocy
Wniosku~\ref{cor:lin7-rzad-transpozycji}.
Każdy radykał jest zerowy
dokładnie wtedy, gdy
\(\operatorname{rk}A=n\),
co jest równoważne
\(\det A\ne0_K\) na mocy
Twierdzenia~\ref{thm:lin6-odwracalnosc}.
Dotyczy to również \(n=0\),
gdy obie podprzestrzenie są zerowe
i \(\det I_0=1_K\).
\end{proof}

\begin{przyklad}[Forma zdegenerowana i niezdegenerowana]
\label{prz:lin8-radykal}
Na \(\mathbb R^2\) forma
\(b((x_1,x_2),(y_1,y_2))
=x_1y_1\)
ma macierz
\(\begin{pmatrix}1&0\\0&0\end{pmatrix}\).
Jej oba radykały są równe
\(\operatorname{span}\{(0,1)\}\),
więc jest zdegenerowana.

Natomiast forma
\(c((x_1,x_2),(y_1,y_2))
=x_1y_2-x_2y_1\)
ma macierz
\(\begin{pmatrix}0&1\\-1&0\end{pmatrix}\)
o wyznaczniku \(1\).
Jest zatem niezdegenerowana.
Równocześnie \(c(v,v)=0\)
dla każdego \(v\in\mathbb R^2\).
Przykład pokazuje, że
naprzemienność i niezdegenerowanie
nie są warunkami sprzecznymi.
\end{przyklad}

\section{Permutowanie argumentów i rodzaje symetrii}
\label{sec:lin8-symetrie}

\begin{definicja}[Formy symetryczne i naprzemienne]
\label{def:lin8-formy-sym-alt}
Niech \(k\geq2\) i
\(V\) będzie przestrzenią nad \(K\).
Forma \(k\)-liniowa
\(F\colon V^k\to K\)
jest \emph{symetryczna}, jeśli
dla dowolnych \(v_1,\ldots,v_k\in V\)
i permutacji \(\sigma\in S_k\)
zachodzi
\[
 F(v_{\sigma(1)},\ldots,v_{\sigma(k)})
 =F(v_1,\ldots,v_k).
\]
Jest \emph{naprzemienna}, jeśli
znika zawsze wtedy, gdy dwa spośród
jej argumentów są równe.
Jest \emph{antysymetryczna}, jeśli
dla wszystkich \(\sigma\in S_k\)
zachodzi
\[
 F(v_{\sigma(1)},\ldots,v_{\sigma(k)})
 =\operatorname{sgn}(\sigma)\,
  F(v_1,\ldots,v_k).
\]
Przestrzenie form symetrycznych
i naprzemiennych oznaczamy odpowiednio
przez \(\operatorname{Sym}^k(V^*)\)
i \(\operatorname{Alt}^k(V^*)\).
Dla \(k=1\) każda forma liniowa
jest symetryczna i naprzemienna;
dla \(k=0\) obie przestrzenie
utożsamiamy z \(K\).
\end{definicja}

\begin{twierdzenie}[Naprzemienność a antysymetria]
\label{thm:lin8-alt-antysym}
Niech \(k\geq2\) i
\(F\colon V^k\to K\)
będzie formą wieloliniową.
Jeśli \(F\) jest naprzemienna,
to jest antysymetryczna.
Jeśli \(2_K\ne0_K\), zachodzi
również implikacja odwrotna.
W charakterystyce \(2\)
odwrotna implikacja może zawodzić.
\end{twierdzenie}
\begin{proof}
Niech \(F\) będzie naprzemienna.
Ustalmy wszystkie argumenty
poza dwoma miejscami, w które
wstawiamy \(u\) i \(v\).
Po wstawieniu w oba miejsca
\(u+v\), naprzemienność daje
\[
 0_K=F(\ldots,u+v,\ldots,u+v,\ldots).
\]
Po rozwinięciu względem obu miejsc
składniki z \((u,u)\) oraz \((v,v)\)
znikają, pozostają więc
\[
 F(\ldots,u,\ldots,v,\ldots)
 +F(\ldots,v,\ldots,u,\ldots)=0_K.
\]
Oznacza to zmianę znaku
przy zamianie dowolnych
dwóch argumentów.
Każda permutacja jest złożeniem
takich zamian, a parzystość
ich liczby wyznacza znak
permutacji (Stwierdzenie~
\ref{prop:lin6-znak-zlozenia}).
Zatem \(F\) jest antysymetryczna.

Odwrotnie, jeśli \(F\)
jest antysymetryczna, to po
zamianie dwóch równych argumentów
jej wartość jest przeciwna
do samej siebie:
\(F(\ldots,v,\ldots,v,\ldots)
=-F(\ldots,v,\ldots,v,\ldots)\).
Jeżeli \(2_K\ne0_K\), dostajemy
\(F(\ldots,v,\ldots,v,\ldots)=0_K\).
Dla charakterystyki \(2\)
wniosek ten nie jest uprawniony.
\end{proof}

\begin{przyklad}[Ostrzeżenie w charakterystyce dwa]
\label{prz:lin8-char2}
Niech \(K=\mathbb F_2\) i \(V=K^2\).
Forma dwuliniowa
\[
 B((x_1,x_2),(y_1,y_2))=x_1y_1
\]
jest symetryczna.
Ponieważ w \(K\) mamy \(-a=a\),
jest także antysymetryczna
w podanym wyżej sensie.
Nie jest jednak naprzemienna,
gdyż \(B((1,0),(1,0))=1_K\).
Nie wolno zatem zastępować
warunku naprzemienności
warunkiem zmiany znaku
bez założenia o charakterystyce.
\end{przyklad}

\begin{stwierdzenie}[Podprzestrzenie form o symetrii]
\label{prop:lin8-sym-alt-podprzestrzenie}
Dla każdej skończenie wymiarowej
przestrzeni \(V\) i każdego
\(k\in\mathbb N_0\)
przestrzenie
\(\operatorname{Sym}^k(V^*)\)
oraz \(\operatorname{Alt}^k(V^*)\)
są podprzestrzeniami
\(\operatorname{Mult}_K(
\underbrace{V,\ldots,V}_{k};K)\).
\end{stwierdzenie}
\begin{proof}
Dla \(k\leq1\) obie przestrzenie
są całą przestrzenią form.
Jeśli \(F,G\) są symetryczne,
to dla każdej permutacji \(\sigma\)
i \(a,b\in K\)
\[
 (aF+bG)(v_{\sigma(1)},\ldots,v_{\sigma(k)})
 =aF(v_1,\ldots,v_k)+bG(v_1,\ldots,v_k).
\]
Zatem \(aF+bG\) jest symetryczna.
Jeśli \(F,G\) są naprzemienne,
to dla krotki z powtórzonym
argumentem obie wartości są zerowe,
więc taka jest również wartość
\(aF+bG\). Każda z obu rodzin
zawiera formę zerową i jest
zamknięta na kombinacje liniowe;
spełnia więc warunek podprzestrzeni.
\end{proof}

\section{Wymiary przestrzeni form symetrycznych}
\label{sec:lin8-wymiar-sym}

\begin{twierdzenie}[Wymiar przestrzeni form symetrycznych]
\label{thm:lin8-wymiar-sym}
Niech \(V\) ma wymiar \(n<\infty\),
a \(k\geq1\).
Wtedy
\[
 \dim\operatorname{Sym}^k(V^*)
 =\binom{n+k-1}{k},
\]
przy konwencji
\(\binom{a}{b}=0\)
dla liczb całkowitych
\(0\leq a<b\).
Dla \(n=0\) i \(k\geq1\)
prawa strona jest równa \(0\).
Dla \(k=0\) wymiar wynosi \(1\),
niezależnie od \(n\).
\end{twierdzenie}
\begin{proof}
Ustalmy uporządkowaną bazę
\(B=(e_1,\ldots,e_n)\) w \(V\).
Jeśli \(F\) jest symetryczna,
to jej wartość na
\((e_{i_1},\ldots,e_{i_k})\)
zależy tylko od multizbioru
indeksów \(i_1,\ldots,i_k\):
dowolna ich permutacja
daje tę samą wartość.
Taką rodzinę współczynników
opisujemy jednoznacznie
przez wartości na krotkach
o indeksach niemalejących:
\[
 1\leq i_1\leq i_2\leq\cdots\leq i_k\leq n.
\]
Odwrotnie, dowolne przyporządkowanie
skalarów tym niemalejącym krotkom
rozszerza się na wszystkie
krotki bazowe przez pominięcie
ich kolejności.
Twierdzenie~\ref{thm:lin8-wspolrzedne}
daje dokładnie jedną formę
wieloliniową o tych wartościach.
Jest ona symetryczna, gdyż
na każdej krotce wektorów bazowych
pozostaje niezmieniona
przy permutowaniu argumentów,
a wieloliniowość rozszerza
tę własność na dowolne wektory.
Przestrzeń form symetrycznych
jest więc izomorficzna z
przestrzenią wszystkich rodzin
skalarów indeksowanych przez
niemalejące krotki.

Gdy \(n\geq1\), podstawienie
\(j_r=i_r+r-1\) daje bijekcję
między takimi krotkami i
ściśle rosnącymi krotkami
\(1\leq j_1<\cdots<j_k\leq n+k-1\).
Jest ich \(\binom{n+k-1}{k}\).
Jeśli \(n=0\), a \(k\geq1\),
nie ma żadnej krotki bazowej,
więc jedyną formą jest zerowa;
wzór również daje zero.
Dla \(k=0\) przestrzeń form
to \(K\), a jej wymiar to \(1\).
\end{proof}

\begin{przyklad}[Formy symetryczne na \(K^2\)]
\label{prz:lin8-sym-baza}
Niech \(V=K^2\), \(k=2\),
z bazą standardową \(e_1,e_2\).
Forma symetryczna \(B\)
jest wyznaczona przez
trzy skalary
\[
 a=B(e_1,e_1),\quad
 b=B(e_1,e_2)=B(e_2,e_1),\quad
 c=B(e_2,e_2).
\]
Współrzędnie
\[
 B((x_1,x_2),(y_1,y_2))
 =ax_1y_1+
 b(x_1y_2+x_2y_1)+cx_2y_2.
\]
Wynika stąd
\(\dim\operatorname{Sym}^2((K^2)^*)=3\),
zgodnie z
\(\binom{2+2-1}{2}=3\).
Rachunek zachowuje ważność
także w charakterystyce \(2\).
\end{przyklad}

\section{Wymiary przestrzeni form naprzemiennych}
\label{sec:lin8-wymiar-alt}

\begin{twierdzenie}[Wymiar i baza form naprzemiennych]
\label{thm:lin8-wymiar-alt}
Niech \(V\) ma wymiar \(n<\infty\),
a \(k\geq1\).
Wtedy
\[
 \dim\operatorname{Alt}^k(V^*)=\binom{n}{k},
\]
gdzie \(\binom{n}{k}=0\) dla \(k>n\).
Dla każdej uporządkowanej bazy
\(B=(e_1,\ldots,e_n)\) przestrzeń
\(\operatorname{Alt}^k(V^*)\)
ma bazę złożoną z form
\(\omega_I\), indeksowanych
ściśle rosnącymi krotkami
\(I=(i_1,\ldots,i_k)\)
o \(1\leq i_1<\cdots<i_k\leq n\),
określonych wzorem
\[
 \omega_I(v_1,\ldots,v_k)
 :=\det\bigl(e^{i_r}(v_s)\bigr)_{r,s=1}^{k},
\]
gdzie \((e^1,\ldots,e^n)\)
jest bazą dualną do \(B\).
Dla \(k=0\) przestrzeń form
naprzemiennych jest równa \(K\)
i ma wymiar \(1=\binom{n}{0}\).
\end{twierdzenie}
\begin{proof}
Niech \(F\) będzie formą naprzemienną.
Jej wartość na krotce wektorów bazowych
z powtórzonym indeksem jest równa zero.
Jeżeli wszystkie indeksy są różne,
to na mocy
Twierdzenia~\ref{thm:lin8-alt-antysym}
wartość na dowolnym przestawieniu
wektorów jest równa wartości
na krotce o indeksach rosnących,
pomnożonej przez znak permutacji.
Zatem forma \(F\) jest jednoznacznie
wyznaczona przez
\(\binom{n}{k}\) skalarów
\[
 F(e_{i_1},\ldots,e_{i_k}),
 \qquad 1\leq i_1<\cdots<i_k\leq n.
\]

Dla takiej krotki \(I\)
macierz \(\bigl(e^{i_r}(v_s)\bigr)\)
zależy liniowo od każdej kolumny
przy ustaleniu pozostałych;
jej wyznacznik jest więc
\(k\)-liniową formą.
Znikanie wyznacznika przy
dwóch identycznych kolumnach
dowodzi jej naprzemienności.
Jeśli \(J=(j_1,\ldots,j_k)\)
jest inną ściśle rosnącą
krotką indeksów, to
\[
 \omega_I(e_{j_1},\ldots,e_{j_k})
 =\det(\delta_{i_rj_s})_{r,s=1}^k.
\]
Dla \(I=J\) macierz jest \(I_k\),
więc wyznacznik wynosi \(1_K\).
Dla \(I\ne J\) istnieje indeks
z \(I\) niewystępujący w \(J\);
odpowiedni wiersz macierzy
jest zerowy, więc wyznacznik
wynosi \(0_K\).
Otrzymujemy więc macierz
wartości \((\omega_I(e_J))_{I,J}\)
równą macierzy jednostkowej.
Z tego wynika liniowa niezależność
form \(\omega_I\).

Każda forma naprzemienna \(F\)
ma te same wartości co
\(\sum_I F(e_{i_1},\ldots,e_{i_k})\omega_I\)
na wszystkich krotkach
o indeksach ściśle rosnących.
Obie formy są naprzemienne,
więc mają identyczne wartości
na wszystkich krotkach bazowych.
Twierdzenie~\ref{thm:lin8-wspolrzedne}
daje ich równość.
Formy \(\omega_I\) stanowią
więc bazę i jest ich
\(\binom{n}{k}\).
Jeśli \(k>n\), każda krotka
\(k\) wektorów bazy zawiera
powtórzenie, więc
jedyna forma naprzemienna
jest zerowa.
Wszystkie przypadki obejmują
także \(n=0\); dla \(k=0\)
obowiązuje wskazana umowa.
\end{proof}

\begin{wniosek}[Wyznacznik jako forma najwyższego stopnia]
\label{cor:lin8-wyznacznik-alt}
Niech \(V\) ma wymiar \(n\geq1\).
Wtedy przestrzeń wszystkich
naprzemiennych form \(n\)-liniowych
na \(V\) ma wymiar \(1\).
Jeśli \(B\) jest uporządkowaną bazą,
to każda taka forma \(F\)
ma jednoznaczną postać
\[
 F(v_1,\ldots,v_n)
 =c\,\det\begin{pmatrix}
 [v_1]_B&\cdots&[v_n]_B
 \end{pmatrix},
 \qquad c\in K.
\]
Jeżeli ponadto
\(F(e_1,\ldots,e_n)=1_K\),
to \(c=1_K\).
\end{wniosek}
\begin{proof}
Z
Twierdzenia~\ref{thm:lin8-wymiar-alt}
wynika
\(\dim\operatorname{Alt}^n(V^*)
=\binom{n}{n}=1\).
Jedyną krotką rosnących
\(n\) indeksów jest
\((1,\ldots,n)\), a zdefiniowana
dla niej forma \(\omega_I\)
ma dokładnie postać
wyznacznika macierzy współrzędnych.
Jej wartość na \(B\)
wynosi \(\det I_n=1_K\).
Pozostałe twierdzenia wynikają
z jednowymiarowości i normalizacji.
\end{proof}

\begin{przyklad}[Naprzemienna forma trójliniowa]
\label{prz:lin8-alt-r3}
Na \(V=\mathbb R^3\) określmy
\(\omega(u,v,w)=\det(u\ v\ w)\),
z kolumnami w bazie standardowej.
Jest to forma trójliniowa i
naprzemienna. Na przykład
\[
 \omega(e_1+e_2,e_2,e_3)=1,
 \qquad
 \omega(e_1,e_2,e_1+e_2)=0.
\]
Pierwszy wynik otrzymujemy,
korzystając z liniowości
w pierwszej kolumnie;
składnik z dwoma kolumnami \(e_2\)
znika. Drugi wyznacznik jest zerowy,
bo trzecia kolumna jest kombinacją
dwóch poprzednich; można również
rozwinąć go według trzeciej kolumny.
Każda inna naprzemienna forma
trójliniowa na \(\mathbb R^3\)
jest wielokrotnością \(\omega\).
\end{przyklad}

\section{Operatory symetryzacji i alternowania}
\label{sec:lin8-operatory-sym-alt}

\begin{definicja}[Nieznormalizowane operatory symetryzacji]
\label{def:lin8-operatory}
Niech \(k\geq1\), a
\(F\colon V^k\to K\)
będzie formą wieloliniową.
Definiujemy \emph{symetryzację}
\(\mathcal S F\) i
\emph{alternowanie}
\(\mathcal A F\) wzorami
\[
 \begin{aligned}
 (\mathcal SF)(v_1,\ldots,v_k)
 &:=\sum_{\sigma\in S_k}
    F(v_{\sigma(1)},\ldots,v_{\sigma(k)}),\\
 (\mathcal AF)(v_1,\ldots,v_k)
 &:=\sum_{\sigma\in S_k}
    \operatorname{sgn}(\sigma)
    F(v_{\sigma(1)},\ldots,v_{\sigma(k)}).
 \end{aligned}
\]
Nie dzielimy przez \(k!\);
dzięki temu obie definicje
mają sens w dowolnej charakterystyce.
\end{definicja}

\begin{twierdzenie}[Własności symetryzacji i alternowania]
\label{thm:lin8-operatory-wlasnosci}
Niech \(k\geq1\).
Operatory \(\mathcal S\) i
\(\mathcal A\) są liniowe.
Dla każdej formy \(k\)-liniowej \(F\)
forma \(\mathcal SF\) jest symetryczna,
a \(\mathcal AF\) jest naprzemienna.
Ponadto:
\[
 \begin{aligned}
 F\text{ symetryczna}&\Longrightarrow
            \mathcal SF=k!\,F,\\
 F\text{ naprzemienna}&\Longrightarrow
            \mathcal AF=k!\,F.
 \end{aligned}
\]
Jeżeli \(k!1_K\ne0_K\),
to operatory
\[
 \mathsf{Sym}:=(k!)^{-1}\mathcal S,
 \qquad
 \mathsf{Alt}:=(k!)^{-1}\mathcal A
\]
są rzutowaniami liniowymi
odpowiednio na przestrzeń form
symetrycznych i naprzemiennych.
\end{twierdzenie}
\begin{proof}
Każdy składnik obu sum jest
wieloliniowy, ponieważ jest
odwzorowaniem \(F\)
ze zmienioną kolejnością argumentów.
Skończona kombinacja takich
odwzorowań jest wieloliniowa.
Liniowość \(\mathcal S\)
i \(\mathcal A\) względem \(F\)
wynika z rozdzielności sum.

Dla dowolnej \(\tau\in S_k\)
podstawienie argumentów
\(v_{\tau(1)},\ldots,v_{\tau(k)}\)
do \(\mathcal SF\) tylko
przestawia składniki sumy po \(S_k\):
złożenie \(\tau\circ\sigma\)
przebiega wszystkie permutacje.
Stąd \(\mathcal SF\) jest symetryczna.

Dla \(\mathcal AF\)
ta sama zmiana indeksów
oraz
\(\operatorname{sgn}(\tau\circ\sigma)
=\operatorname{sgn}(\tau)
 \operatorname{sgn}(\sigma)\)
dają
\[
 (\mathcal AF)(v_{\tau(1)},\ldots,v_{\tau(k)})
 =\operatorname{sgn}(\tau)
  (\mathcal AF)(v_1,\ldots,v_k).
\]
Aby wykazać \emph{naprzemienność}
bez założenia o charakterystyce,
załóżmy, że argumenty na
pozycjach \(p\ne q\) są równe.
Permutacje sumy alternującej
parujemy przez
\(\sigma\mapsto(p\,q)\circ\sigma\).
Dla każdej pozycji \(j\) mamy
\[
 v_{((p\,q)\circ\sigma)(j)}
 =v_{\sigma(j)},
\]
ponieważ \(v_p=v_q\):
zamiana wartości indeksów \(p\) i \(q\)
nie zmienia żadnego z argumentów.
Zatem składniki odpowiadające
\(\sigma\) oraz \((p\,q)\circ\sigma\)
mają tę samą wartość \(F\).
Ich znaki są przeciwne na mocy
Stwierdzenia~\ref{prop:lin6-znak-zlozenia},
więc suma takiej pary jest \(0_K\),
również w charakterystyce \(2\).
Odwzorowanie
\(\sigma\mapsto(p\,q)\circ\sigma\)
jest inwolucją bez punktów stałych,
gdyż \(p\ne q\).
Zatem \(\mathcal AF\)
znika przy dwóch równych argumentach,
również w charakterystyce \(2\).

Jeśli \(F\) jest symetryczna,
to wszystkie \(k!\) składników
\(\mathcal SF\) są równe \(F\);
zatem \(\mathcal SF=k!F\).
Jeśli \(F\) jest naprzemienna,
to każdy składnik \(\mathcal AF\)
jest równy
\(\operatorname{sgn}(\sigma)^2 F=F\).
Stąd \(\mathcal AF=k!F\).

Gdy \(k!1_K\ne0_K\), w ciele
możemy dzielić przez \(k!\).
Obrazy \(\mathsf{Sym}\) i
\(\mathsf{Alt}\) zawierają się
w odpowiednich podprzestrzeniach,
a na każdej z nich operator
jest identycznością.
Dlatego oba są rzutowaniami:
\(\mathsf{Sym}^2=\mathsf{Sym}\)
i \(\mathsf{Alt}^2=\mathsf{Alt}\).
\end{proof}

\begin{uwaga}[Co dzieje się przy dzieleniu przez \(k!\)]
\label{rem:lin8-char-faktorial}
Warunek \(k!1_K\ne0_K\)
jest automatycznie spełniony
w charakterystyce \(0\), a w
dodatniej charakterystyce \(p\)
dokładnie wtedy, gdy \(p>k\).
Jeśli \(k!1_K=0_K\), to
operator \(\mathcal S\)
zeruje każdą formę już symetryczną,
a operator \(\mathcal A\)
zeruje każdą formę już naprzemienną.
Nie wolno wtedy dzielić przez \(k!\)
ani wnioskować, że znormalizowane
rzutowania istnieją na podstawie
powyższych wzorów.
Mimo to same podprzestrzenie
form symetrycznych i naprzemiennych
oraz wzory na ich wymiary pozostają
poprawne.
\end{uwaga}

\begin{przyklad}[Alternowanie nie jest zawsze rzutowaniem]
\label{prz:lin8-alt-f2}
Na \(V=\mathbb F_2^2\) niech
\(F(x,y)=x_1y_2\).
Wtedy
\[
 (\mathcal AF)(x,y)
 =x_1y_2-x_2y_1
 =x_1y_2+x_2y_1,
\]
co jest niezerową formą naprzemienną.
Jednak \(\mathcal A(\mathcal AF)=
2\,\mathcal AF=0\)
w \(\mathbb F_2\), mimo że
\(\mathcal AF\ne0\).
Zatem sam operator
\(\mathcal A\) nie jest
w tym przykładzie rzutowaniem.
\end{przyklad}

\section{Wartości na przekątnej i polaryzacja}
\label{sec:lin8-polaryzacja}

\begin{definicja}[Funkcja diagonalna formy]
\label{def:lin8-diagonalna}
Niech \(k\geq1\),
\(F\colon V^k\to K\)
będzie formą \(k\)-liniową.
Jej \emph{funkcją diagonalną}
nazywamy
\[
 p_F\colon V\to K,\qquad
 p_F(v):=F(v,\ldots,v).
\]
Funkcja \(p_F\) spełnia
\(p_F(av)=a^kp_F(v)\)
dla wszystkich \(a\in K\)
i \(v\in V\).
Na ogół nie jest liniowa.
\end{definicja}

\begin{twierdzenie}[Wzór polaryzacyjny]
\label{thm:lin8-polaryzacja}
Niech \(k\geq1\),
\(V\) będzie przestrzenią nad \(K\),
a \(F\colon V^k\to K\)
\emph{symetryczną} formą \(k\)-liniową.
Dla dowolnych \(v_1,\ldots,v_k\in V\)
zachodzi
\[
 \sum_{S\subset\{1,\ldots,k\}}
 (-1)^{k-|S|}
 p_F\!\left(\sum_{i\in S}v_i\right)
 =k!\,F(v_1,\ldots,v_k).
\]
Suma po pustym \(S\)
jest poprawnie określona przez
\(p_F(0_V)=0_K\).
Jeśli \(k!1_K\ne0_K\),
funkcja \(p_F\) wyznacza \(F\)
jednoznacznie przez podzielenie
prawej strony wzoru przez \(k!\).
\end{twierdzenie}
\begin{proof}
Dla ustalonego \(S\)
rozwijamy funkcję diagonalną:
\[
 p_F\!\left(\sum_{i\in S}v_i\right)
 =\sum_{(j_1,\ldots,j_k)\in S^k}
   F(v_{j_1},\ldots,v_{j_k}).
\]
Zamieńmy kolejność
skończonych sum po \(S\)
i po krotkach indeksów.
Dla konkretnej krotki
\((j_1,\ldots,j_k)\) niech
\(J=\{j_1,\ldots,j_k\}\).
Jej wartość \(F(v_{j_1},\ldots,v_{j_k})\)
pojawia się dokładnie w tych
sumach, dla których \(J\subset S\),
z łącznym współczynnikiem
\[
 \sum_{J\subset S\subset\{1,\ldots,k\}}
  (-1)^{k-|S|}
 =(1-1)^{k-|J|}.
\]
Ostatni zapis jest tylko
skrótem dwumianowego rachunku
sumy; gdy \(|J|<k\),
współczynnik wynosi \(0_K\).
Gdy \(|J|=k\), wynosi \(1_K\).
Pozostają więc dokładnie
te krotki \((j_1,\ldots,j_k)\),
w których każdy z indeksów
\(1,\ldots,k\) występuje raz,
czyli \(k!\) krotek opisanych
permutacjami. Ponieważ \(F\)
jest symetryczna, każda daje
wartość \(F(v_1,\ldots,v_k)\).
Sumując, otrzymujemy żądaną równość.
Jeśli \(k!1_K\ne0_K\),
jest to element odwracalny w ciele,
co pozwala rozwiązać wzór
względem \(F(v_1,\ldots,v_k)\).
\end{proof}

\begin{wniosek}[Polaryzacja formy dwuliniowej]
\label{cor:lin8-polaryzacja-dwu}
Jeśli \(2_K\ne0_K\)
i \(b\colon V\times V\to K\)
jest symetryczna, to
dla \(q(v)=b(v,v)\) zachodzi
\[
 b(u,v)=\frac{q(u+v)-q(u)-q(v)}{2_K}.
\]
Dla dowolnej charakterystyki
prawdziwa jest bez dzielenia
tożsamość
\(q(u+v)-q(u)-q(v)=2_Kb(u,v)\).
\end{wniosek}
\begin{proof}
Z dwuliniowości otrzymujemy
\[
 \begin{aligned}
 q(u+v)
 &=b(u+v,u+v)\\
 &=b(u,u)+b(u,v)+b(v,u)+b(v,v).
 \end{aligned}
\]
Symetria daje
\(b(v,u)=b(u,v)\),
więc różnica z tezy wynosi
\(2_Kb(u,v)\).
Jeśli \(2_K\ne0_K\), dzielimy
przez tę niezerową liczbę.
\end{proof}

\begin{przyklad}[Funkcja diagonalna nie zawsze odtwarza formę]
\label{prz:lin8-diagonalna-char2}
Niech \(K=\mathbb F_2\),
\(V=K^2\), a
\[
 b(x,y)=x_1y_2+x_2y_1.
\]
Forma \(b\) jest symetryczna,
niezerowa i naprzemienna.
Dla każdego \(x=(x_1,x_2)\)
mamy
\[
 q(x)=b(x,x)=2_Kx_1x_2=0_K.
\]
Zatem \(b\) i forma zerowa
mają tę samą funkcję diagonalną,
mimo że \(b(e_1,e_2)=1_K\).
Jest to dokładny powód,
dla którego w twierdzeniu
o polaryzacji wymagamy
odwracalności \(k!\).
\end{przyklad}

\section{Podsumowanie}
\label{sec:lin8-podsumowanie}

Odwzorowanie \(k\)-liniowe
jest liniowe w każdym argumencie,
lecz na ogół nie jest liniowe
na produkcie przestrzeni
rozumianym jako jedna
przestrzeń wektorowa.
Jego wartości na krotkach
wektorów bazowych wyznaczają
je całkowicie i pozwalają
obliczać wymiar przestrzeni
odwzorowań wieloliniowych.
Formy symetryczne i naprzemienne
tworzą podprzestrzenie
o wymiarach odpowiednio
\(\binom{n+k-1}{k}\) i
\(\binom{n}{k}\).
Drugi wzór prowadzi do
jednowymiarowej przestrzeni
form naprzemiennych
\(n\)-liniowych na \(K^n\),
a więc do naturalnego ujęcia
wyznacznika.
Operatory symetryzacji i alternowania
są poprawnie zdefiniowane
w każdym ciele, ale ich normalizacja
przez \(k!\) wymaga, aby
\(k!\ne0_K\).
Z tej samej przyczyny odtwarzanie
formy symetrycznej z jej funkcji
diagonalnej wymaga odpowiedniego
warunku na charakterystykę.
